\section{Introduction}
\label{sec:introduction}

Public procurement records describe how public buyers solicit work, choose suppliers, and publish
awards and contracts.  The Open Contracting Data Standard (OCDS) gives these disclosures a common
model covering parties, tenders, awards, contracts, and linked documents \citep{ocds-standard}.
Data published this way can answer questions such as which suppliers serve several buyers,
whether disclosure is complete for a class of notices, or how activity shifts across years.  None
of these is a document lookup.  An answer may depend on the role an organisation plays in one
record, on a set of organisations derived from that record, on a second retrieval bound by the
derived set, and on a final count, comparison, or ranking.

Finding \emph{relevant} evidence is not the same as recovering the \emph{complete} denotation an
aggregate needs.  Let $R(q)$ be the full set of contract records satisfying question $q$, and let
a top-$k$ retriever return $\widehat{R}(q)\subset R(q)$.  Even when $\widehat{R}(q)$ holds
convincing examples, it does not follow that
\begin{equation}
 |\widehat{R}(q)|=|R(q)|, \qquad
 \sum_{r\in\widehat{R}(q)}v(r)=\sum_{r\in R(q)}v(r),
 \label{eq:retrieval-not-aggregation}
\end{equation}
nor that a ranking computed from the retrieved subset shares the true maximum.  Retrieval-centred
KGQA systems find answer-bearing subgraphs well \citep{sun-etal-2019-pullnet, jiang2023unikgqa,
luo2024rog}.  Exact counts and set operations need something further: an explicit scope, typed
intermediate values, and exhaustive execution inside a declared data universe.

Letting a language model write the answer directly fails differently.  A model can produce a
well-formed query that swaps buyer and supplier roles, projects a contract identifier where the
question asked for an organisation name, or sums a non-additive monetary field.  The answers stay
fluent and look plausible.  Prior work has responded by representing complex questions as
executable programs \citep{cao-etal-2022-kqa, wolfson-etal-2020-break}, constraining formal
generation \citep{scholak-etal-2021-picard, shu-etal-2022-tiara}, and feeding execution or
environment signals back into the parser \citep{chen-etal-2021-retrack,
huang-etal-2024-queryagent}.  We take that starting point and push it into a division of labour:
the model proposes a typed computation, and the data system decides every factual value.

We build this at two levels, and keep them separate in the evaluation as well as in the code.  A
compact model is trained to translate procurement questions into a closed algebra of record
filters, distinct projections, semi- and anti-joins, set operations, grouping, aggregation,
comparison, and ranking; each program it emits is type checked and executed over a frozen record
universe, and an offline scorer compares the result with a fixed oracle.  Around that model, a
deployment runtime adds a structured understanding stage and a conditional planner.  Intents that
compile deterministically take a compiler path and the rest go to a typed graph planner, after
which the plan is grounded against allowed schema fields, executed exhaustively, and put through
preflight, postflight, answer-sanity, and evidence checks.  Structural defects can trigger
bounded replanning, but a clean empty result ends the process, so the system cannot keep searching
until it finds an answer it prefers.  The first protocol never invokes the understander, planner,
evidence bundle, or repair loop, which stops an offline oracle check from being reported as an
inference-time verifier and stops a component that was never called from taking credit for an
item.

\begin{figure}[t]
  \centering
  \includegraphics[width=\linewidth]{figures/fig01_running_example.png}
  \caption{One bridge question over the frozen procurement graph.  Retrieval-augmented reading
  retrieves 40 records, among them all 23 naming the supplier, and answers 14; an untuned model
  of the same size emits a program that type checks and executes cleanly but stops after the
  first hop and returns 23.  Our runtime executes the two-hop bridge exhaustively---23 source
  records, 16 buyer values, 4{,}257 notices---and releases an answer only once execution,
  evidence, and answer-sanity checks pass.  Every value comes from one saved run; its provenance
  and limits are given in Appendix \ref{app:examples}.}
  \label{fig:visual-thesis}
\end{figure}

The experiments use a provenance-preserving graph built from UK procurement releases for
2022--2026 and a frozen \PacsN{}-question compositional test.  The typed-program model reaches
\PacsOursA\% strict accuracy on the canonical surface, \PacsGainBase{} points above the same
untuned 8B base model and \PacsGainTeacher{} points above a cloud planner, and loses only
\PacsSurfaceDrop{} points under independent surface naturalisation.  Gains concentrate on
overlap/exclusion and relational composition, while supplier concentration and
disclosure/compliance stay much harder.  A separate balanced \RuntimeN{}-item artifact puts the
complete runtime at \RuntimeOurs\% against \RuntimeBase\% untuned; we treat it as secondary
evidence, because its stored outputs travel less well than the PACS artifacts.

Four things follow from this design.  The first is the procurement graph itself, with a
deterministic query-record projection and explicit entity-resolution, additivity, and coverage
checks, which is what makes any denotation over it well defined.  The second is the closed typed
algebra together with the program-first supervision pipeline that teaches a compact model to
write in it.  The third is a runtime that keeps learned understanding and planning apart from
grounding, exhaustive execution, evidence checks, bounded repair, and controlled non-answering.
The fourth is PACS, a frozen compositional test with two surface channels, operator-family and
depth diagnostics, answerability controls, and a published account of where its isolation and
scoring procedures fall short.  None of this guarantees that the source data are complete in the
world, that stored organisation names are unique legal entities, or that passing an execution
check proves semantic correctness.  The claim we do make is testable: inside a declared frozen
data universe, the computation a model proposes can be type checked, executed, traced, and
rejected when its implemented preconditions fail.
